El vínculo electrodébil parte del mismo operador angular que produjo las constitutivas del vacío, pero utiliza otro funcional.\par

Para obtener \(\mu\), \(\varepsilon\), \(Z\) y \(c\) había que conservar el espectro completo. Había que distinguir las dos hojas.\par

El determinante hace lo contrario:\par

\[
D_A=\det T=q_+q_-.
\]

Al multiplicar los autovalores, el determinante integra las dos orientaciones en una magnitud angular común e invariante. Publica la parte par y reserva la identidad de hoja en la memoria orientada.\par

Por eso \(D_A\) es par bajo la inversión\par

\[
C^*\longmapsto-C^*.
\]

El sector \(W/Z\) posee justamente esa estructura. Las rutas de \(W\) y \(Z\) son adyacentes en la coordenada angular principal y comparten la misma coordenada orientada. Al formar el cociente, la contribución de \(C^*\) se cancela.\par

Queda\par

\[
\left(\frac{m_W}{m_Z}\right)^2
=
D_A,
\]

y por tanto\par

\[
\sin^2\theta_W
=
1-D_A.
\]

La identidad revela que el ángulo de Weinberg es el defecto complementario del determinante angular.\par

El determinante mide la parte común que sobrevive al colapsar las dos hojas en su producto. \(1-D_A\) mide la parte complementaria necesaria para la mezcla.\par

Podríamos decirlo así: el vacío conserva toda la respuesta; el sector electrodébil utiliza la compresión par de esa respuesta.\par

El mismo antecedente angular produce dos funciones distintas:\par

\begin{itemize}[leftmargin=*,itemsep=0.35em]
\item el cálculo funcional completo genera las relaciones constitutivas;
\item el determinante genera el invariante de mezcla.
\end{itemize}

Eso explica por qué vacío y ángulo de Weinberg son dos realizaciones funcionales distintas de un mismo antecedente angular.\par

A partir de ahí, los acoplamientos \(g\) y \(g'\) son, en el nivel declarado, las dos proyecciones del mismo acoplamiento electromagnético racionalizado sobre las direcciones seno y coseno de la mezcla:\par

\[
g^2
=
\frac{4\pi\alpha_{\rm HMT}}{1-D_A},
\qquad
g'^2
=
\frac{4\pi\alpha_{\rm HMT}}{D_A}.
\]

La identidad custodial de árbol\par

\[
\rho_{\rm EW}^{(0)}=1
\]

expresa que la misma partición angular organiza coherentemente masa y acoplamiento.\par

La constante de Fermi aparece después como la respuesta efectiva de baja energía del canal cargado:\par

\[
G_F^{(0)}
=
\frac{\pi\alpha_{\rm HMT}}
{\sqrt2(1-D_A)m_W^2}.
\]

La constante de Fermi aparece como una composición de la estructura fina, el defecto angular y la escala del portador \(W\).\par

El sector de Higgs introduce una diferencia decisiva. Su ruta posee una coordenada orientada distinta de la ruta de \(W\), por lo que \(C^*\) persiste. El autoacoplamiento escalar publica la orientación que el determinante reserva.\par

El sector de calibre es par: utiliza \(D_A\).\par
El sector escalar es orientado: retiene \(C^*\).\par

Esta división resulta físicamente sugerente. Los campos de calibre organizan la parte común de la transformación; el sector escalar conserva la memoria de la orientación rota.\par

Las correcciones radiativas pertenecen a un nivel posterior y deben construirse prospectivamente desde los bucles, estados y escalas internos, con \(G_F\) como salida y evaluación terminal. El esqueleto causal ya está presente: vacío, mezcla, acoplamientos y respuesta de Fermi proceden de distintas lecturas del mismo estado angular enriquecido.\par

